\lecture{1}{Wed 13 Sep 2023 10:43}{Upcrossing Inequality}

\begin{proof}
	Define stopping times
	\begin{align*}
		N_0 &=  -1 \\
		N_1 &= \inf \{m > N_0: X_m \le a\}  \\
		N_2 &= \inf \{m > N_1: X_m \ge b\}
	.\end{align*}
	In general, $N_{2k-1}$ are downcrossing times of $a$, $N_{2k}$ are upcrossing times of $b$.

	Then
	\[
		H_m = \begin{cases}
			1 & N_{2k-1} < m \le N_{2k} \text{ for some $k$}\\
			0 
		\end{cases}
	.\] 
	is a predictable sequence. 

	\begin{intuition}
		$H$ is a ``switch'' that turns on during an upcrossing and turns off during a downcrossing.
	\end{intuition}

	Now define
	\[
		Y_{n} = (X_n - a)^{+} + a = \begin{cases}
			a & X_n < a \\
			X_n & X_n \ge a
		\end{cases}
	.\] 
	We see that $Y_n$ is a submartingale.

	Let $(H * Y)_y = \sum_{k - 1}^n H_n (Y_n - Y_{n-1})$ be the martingale transformation of $Y$ by $H$. Then $(H * Y)_N$ corresponds to the sum of variation of $Y$ across each upcrossing. For each complete upcrossing, the difference is at least $b-a$ ; it is possible that the last upcrossing is incomplete by time $N$, but the contribution is nonnegative in expectation since we are working with a submartingale.
	\[
		(b-a) U_N[a,b] \le (H * Y)_N
	.\] 

	Then we have $Y_n - Y_0 = (H * Y)_n + ((1-H)*Y)_n$. But $1 - H_n \ge 0$, so 
	\[
		E((H*Y)_N) =  E(Y_N - Y_0) - E\left( ((1-H)*Y)_N \right)  \le E(Y_N - Y_0)
	.\] 
\end{proof}

\begin{remark}
	This is a standard argument involving martingales. We count the upcrossings, and replace $X_n$ with $Y_n$ to make it a submartingale.
\end{remark}

\begin{corollary}
	If $X_n \ge 0$ is a supermartingale, then $X_n \to X_{\infty }$ a.s.
\end{corollary}
\begin{proof}
	Apply submartingale theorem to $Y_n = -X_n \le 0$. Then
	$\sup _n E(Y_n^+) < \infty $. Then $Y_n \to Y < \infty $ a.s., so $X_n \to  X  > - \infty $ a.s.
	Furthermore, 
	Fatou's Lemma gives 
	\[
		EX \le \lim\inf  EX_n \le EX_0
	.\] 
\end{proof}

\begin{theorem}[Borel-Cantelli Lemma]
	Take $\mathcal{F}_n$ a filtration, $\mathcal{F}_0 = \{\emptyset , \Omega\} $. Then if $B_n \in \mathcal{F}_n$, we have
	\[
		\{B_n \text{ i.o.}\} \equiv \left\{ \sum_n 1_{B_n} = \infty  \right\}  = \left\{\sum_n P\left(B_n | \mathcal{F}_{n-1}\right)  = \infty \right\} 
	.\] 

	Furthermore, they go to infinity at the same rate:
	\[
		\lim_{n \to \infty }
		\frac{\sum_i^n B_i}{\sum_i^n P\left( B_n | \mathcal{F}_{n-1} \right) }
	.\] 
\end{theorem}



